\newcommand{\kB}{\mathcolor{gray}{B}}
\newcommand{\kC}{\mathcolor{Red}{C}}
\newcommand{\kD}{\mathcolor{Blue}{D}}
\newcommand{\kE}{\mathcolor{YellowOrange}{E}}
\newcommand{\kF}{\mathcolor{Magenta}{F}}

\section{Zadanie 8}
\autor{Dominik Korsa}
Dostaliśmy od Świętego Mikołaja \(\kB\) białych kul, gdzie \(\kB\) jest zmienną
o rozkładzie Poissona z parametrem \(\mu = 2023\).

Każda kula z prawdopodobieństwem \(p\) (\(0 < p < 1\)) jest
\textcolor{Red}{pomalowana na czerwono} niezależnie od innych kul.
Niech \(\kC\) będzie zmienną losową oznaczającą liczbę czerwonych kul w tym
momencie.


Każda \textcolor{Red}{czerwona kula} z prawdopodobieństwem \(p\) jest
\textcolor{Blue}{pomalowana na niebiesko}
niezależnie od innych kul. Niech \(\kD\) będzie zmienną losową oznaczającą liczbę
niebieskich kul w tym momencie.

Każda \textcolor{Blue}{niebieska kula} z prawdopodobieństwem \(p\) jest
\textcolor{YellowOrange}{pomalowana na żółto}
niezależnie od innych kul. Niech \(\kE\) będzie zmienną losową oznaczającą liczbę
żółtych kul w tym momencie.

\begin{enumerate}
	\item\label{set06.task08.1}
	      Dla \(p = \sqrt[3]{1 - \frac{\ln 3}{2023}}\) sprawdź, czy wartość
	      \(P(\kB > \kE)\) jest większa, mniejsza, czy równa \(\frac{2}{3}\).

	\item\label{set06.task08.2}
	      Dla jakiego \(p\) wartość oczekiwana \(\kE\) jest równa \(1\)?

	\item\label{set06.task08.3}
	      Znajdź \(\ev{\kE \mid \kB}\) i \(\ev{\kB \mid \kE}\).
\end{enumerate}

\subsection*{Rozwiązanie podpunktu \eqref{set06.task08.1}}
\(P(\kB > \kE)\) to prawdopodobieństwo zdarzenia, że \textbf{nie wszystkie} kule
zostały \textcolor{YellowOrange}{pomalowane na żółto} (a więc wcześniej też
\textcolor{Red}{na czerwono} i \textcolor{Blue}{na niebiesko}).

Policzmy prawdopodobieństwo zdarzenia przeciwnego, czyli \(P(\kB = \kE)\).

\begin{align*}
	P(\kB = \kE)
	 & = \sum_{i = 0}^\infty P(\kB = \kE \cap \kB = i)                  \\
	 & = \sum_{i = 0}^\infty P(\kB = \kE \mid \kB = i) \cdot P(\kB = i) \\
\end{align*}

W sumowanym wyrażeniu \(P(\kB = \kE \mid \kB = i)\) to prawdopodobieństwo zdarzenia,
że wszystkie z \(i\) wylosowanych kul zostały pokolorowane
\textcolor{Red}{na czerwono}, \textcolor{Blue}{na niebiesko}
i \textcolor{YellowOrange}{na żółto}, więc:
\[ P(\kB = \kE \mid \kB = i) = p^{3i}. \]

\(B\) ma rozkład Poissona, więc
\(\displaystyle P(\kB = i) = \frac{e^{-\mu} \cdot \mu^i}{i!}\).

\needspace{10ex}
Podstawmy te wartości w powyższej równości i przekształćmy ją tak, aby po prawej
stronie znaleźć prawdopodobieństwo wartości w rozkładzie Poissona z parametrem
\(p^3 \cdot \mu\):
\begin{align*}
	P(\kB = \kE)
	 & = \sum_{i = 0}^\infty p^{3i} \cdot \frac{e^{-\mu} \cdot \mu^i}{i!}                                                     \\
	 & = \sum_{i = 0}^\infty \frac{e^{-\mu} \cdot (p^3 \cdot \mu)^i}{i!}                                                      \\
	 & = \sum_{i = 0}^\infty \frac{e^{-\mu} \cdot e^{(p^3 \cdot \mu)} \cdot e^{-(p^3 \cdot \mu)} \cdot (p^3 \cdot \mu)^i}{i!} \\
	 & = e^{(p^3 - 1)\mu} \cdot \sum_{i = 0}^\infty \frac{e^{-(p^3 \cdot \mu)} \cdot (p^3 \cdot \mu)^i}{i!}                   \\
\end{align*}

Powiedzmy, że \(X\) to zmienna losowa o rozkładzie Poissona z parametrem
\(p^3 \cdot \mu\). Wtedy:
\[
	\sum_{i = 0}^\infty \frac{e^{-(p^3 \cdot \mu)} \cdot (p^3 \cdot \mu)^i}{i!}
	= \sum_{i = 0}^\infty P(X = i) = 1.
\]
Mamy więc, że \(P(\kB = \kE) = e^{(p^3-1)\mu}\).
Podstawiając wartości z zadania, czyli \(p = \sqrt[3]{1 - \frac{\ln 3}{2023}}\)
i \(\mu = 2023\), mamy:
\[
	P(\kB = \kE)
	= e^{\Big(\left(\sqrt[3]{1 - \frac{\ln 3}{2023}}\right)^3-1\Big) \cdot 2023}
	= e^{-\frac{\ln 3}{2023} \cdot 2023}
	= e^{- \ln 3}
	= \frac{1}{3},
\]
zatem \(\displaystyle
P(\kB > \kE) = 1 - P(\kB = \kE) = \frac{2}{3}
\).

\subsection*{Rozwiązanie podpunktu \eqref{set06.task08.2}}
Intuicyjnym jest, że oczekiwana liczba kul pokolorowanych
\textcolor{YellowOrange}{na żółto}
będzie iloczynem oczekiwanej liczby wszystkich kul i prawdopodobieństwa tego,
że dana kula zostanie pokolorowana na żółto, czyli:
\(\ev{\kE} = p^3 \cdot \ev{\kB}\).
Poniższy fragment jedynie formalizuje tę intuicję:

\begin{quote}
	Niech \(\kE_j\) to indykator zdarzenia, że \(j\)-ta kula została pomalowana
	\textcolor{YellowOrange}{na żółto}:
	\[
		E[\kE_j] = P(\kE_j = 1) = p^3.
	\]
	W zależności od wartości zmiennej losowej \(\kB\):
	\[
		\kE = \sum_{j = 1}^{\kB} \kE_j.
	\]

	Wartość oczekiwana \(E\) to
	\(\displaystyle
	\ev{\kE} = \sum_{i = 0}^{\infty}
	\Big( \ev{\kE \mid \kB = i} \cdot P(\kB = i) \Big)
	\).\\
	Obliczmy \(\ev{\kE \mid \kB = i}\):
	\[
		\ev{\kE \mid \kB = i}
		= \ev{\sum_{j = 1}^{\kB} \kE_j \mid \kB = i}
		% &= \ev{\sum_{j = 0}^{i} \kE_j}\\
		= \sum_{j = 1}^{i} \ev{\kE_j}
		= i \cdot p^3
	\]

	Podstawiając tę wartość do \(\ev{\kE}\):
	\begin{align*}
		\ev{\kE}
		 & = \sum_{i = 0}^{\infty}
		\Big( \ev{\kE \mid \kB = i} \cdot P(\kB = i) \Big)                  \\
		 & = \sum_{i = 0}^{\infty} \Big( i \cdot p^3 \cdot P(\kB = i) \Big) \\
		% &= p^3 \cdot \sum_{i = 0}^{\infty} \Big( i \cdot P(\kB = i) \Big)\\
		 & = p^3 \cdot \ev{\kB}                                             \\
	\end{align*}
\end{quote}

\(\kB\) jest zmienną losową z rozkładem Poissona, więc jej wartość oczekiwana to
\(\ev{\kB} = \mu = 2023\).

Aby \(\ev{\kE} = 1\):
\(\displaystyle
p^3 = \frac{\ev{\kE}}{\ev{\kB}}
= \frac{1}{2023}
\),
czyli \(\displaystyle
p = \frac{1}{\sqrt[3]{2023}}
\).

\subsection*{Rozwiązanie podpunktu \eqref{set06.task08.3}}
\subsubsection*{Część pierwsza: \(\ev{\kE \mid \kB}\)}
\(\ev{\kE \mid \kB}\) jest funkcją
\(\ev{\kE \mid \kB}\colon \mathbb{R} \to \mathbb{R}\) zdefiniowaną jako:
\[
	\ev{\kE \mid \kB}(i) := \ev{\kE \mid \kB = i}.
\]
Wartość \(\ev{\kE \mid \kB = i}\) policzyliśmy już w podpunkcie \eqref{set06.task08.2}.\\
Dla \(i \in \mathbb{N}\): \(\ev{\kE \mid \kB}(i) = i \cdot p^3\),
dla pozostałych \(i\) \(\ev{\kE \mid \kB}(i) = 0\), bo \(P(\kB = i) = 0\).

\subsubsection*{Część druga: \(\ev{\kB \mid \kE}\)}
\(\ev{\kB \mid \kE = i}\) to oczekiwana liczba wszystkich kul zakładając,
że \(i\) spośród tych kul zostało pokolorowanych \textcolor{YellowOrange}{na żółto}.

Mogłoby się wydawać, że \(\ev{\kB \mid \kE = i} = i \cdot p^{-3}\),
ale tak nie jest, bo \(\ev{\kB \mid \kE = 0} > 0\).

Przeprowadźmy inne doświadczenie losowe.
Losujemy niezależnie:
\begin{itemize}
	\item Pewną liczbę \(\kE\) żółtych kul, gdzie \(\kE\) ma rozkład Poissona z
	      parametrem \(\mu_{\kE} = \mu \cdot p^3\).
	\item Pewną liczbę \(\kF\) kul w innych kolorach, gdzie \(\kF\) ma rozkład
	      Poissona z parametrem \(\mu_{\kF} = \mu \cdot (1-p^3)\).
\end{itemize}
Zmienna losowa \(\kB\) to liczba wszystkich kul, więc \(\kB = \kE + \kF\).

\textbf{Lematu z wykładu}: Suma zmiennych o rozkładzie Poissona z parametrami \(\mu_1\)
i \(\mu_2\) jest zmienną o rozkładzie Poissona z parametrem \(\mu_1 + \mu_2\).

\textbf{Blef alert}: ,,Utożsamianie zmiennych'' poniżej wymaga lepszego dowodu.
Używamy tutaj związku zmiennych o rozkładzie Poissona ze zmiennymi o rozkładzie
dwumianowym, ale ja ten związek niezbyt rozumiem.

Na mocy tego lematu \(\kB\) jest zmienną o rozkładzie Poissona z parametrem
\(\mu_{\kE} + \mu_{\kF} = \mu\), co zgadza się z \(\kB\) z treści zadania
--- utożsammy te zmienne ze sobą.

Dla ustalonego \(\kB\) zmienna \(\kE\) ma rozkład dwumianowy z parametrem \(p^3\),
co również jest zgodne z rozkładem \(\kE\) z treści zadania, więc je także możemy
utożsamić.

\begin{align*}
	\ev{\kB \mid \kE = i}
	 & = \ev{(\kE + \kF) \mid \kE = i}                 \\
	 & = \ev{\kE \mid \kE = i} + \ev{\kF \mid \kE = i} \\
	 & = i + \ev{\kF \mid \kE = i}                     \\
\end{align*}
Zmienne \(\kE\) i \(\kF\) są niezależne, więc
\(\ev{\kF \mid \kE = i} = \ev{\kF} = \mu_{\kF} = \mu \cdot (1 - p^3)\).

\[
	\ev{\kB \mid \kE = i} = i + \mu (1-p^3)
	= i + 2023 \cdot (1-p^3).
\]
Ponownie z definicji \(\ev{\kB \mid \kE}(i) = \ev{\kB \mid \kE = i}\)
dostajemy odpowiedź.

\subsection*{Odpowiedź}
Podpunkt \eqref{set06.task08.1}: \(\boxed{P(\kB > \kE) = \frac{2}{3}}\).\\
Podpunkt \eqref{set06.task08.2}: \(\boxed{p = \frac{1}{\sqrt[3]{2023}}}\)\\
Podpunkt \eqref{set06.task08.3}: \(\boxed{
	\ev{\kE \mid \kB}(\kB) = \begin{cases}
		\kB \cdot p^3 & \text{gdy } \kB \in \mathbb{N}, \\
		0             & \text{wpp.}
	\end{cases}
}\)\\
oraz \(\boxed{
	\ev{\kB \mid \kE}(\kE) = \begin{cases}
		\kE + 2023 \cdot (1-p^3) & \text{gdy } \kE \in \mathbb{N}, \\
		0                        & \text{wpp.}
	\end{cases}
}\)

\undef\kB
\undef\kC
\undef\kD
\undef\kE
